%
\vspace{-0.1cm}
\section{Reinforcement Learning Background}
\label{sec:background}

We consider the standard reinforcement learning setting where an agent interacts with an environment $\mathcal{E}$ over a number of discrete time steps.
%
At each time step $t$, the agent receives a state $s_t$ and selects an action $a_t$ from some set of possible actions $\mathcal{A}$ according to its policy $\pi$, where $\pi$ is a mapping from states $s_t$ to actions $a_t$.
In return, the agent receives the next state $s_{t+1}$ and receives a scalar reward $r_t$.
The process continues until the agent reaches a terminal state after which the process restarts.
%
The return $R_t = \sum_{k=0}^{\infty} \gamma^k r_{t+k}$ is the total accumulated return from time step $t$ with discount factor $\gamma \in (0,1]$.
The goal of the agent is to maximize the expected return from each state $s_t$.

The action value $Q^{\pi}(s,a) = \mathbb{E}\left[R_t|s_t=s, a\right]$ is the expected return for selecting action $a$ in state $s$ and following policy $\pi$.
The optimal value function $Q^*(s,a) = \max_{\pi} Q^{\pi}(s,a)$ gives the maximum action value for state $s$ and action $a$ achievable by any policy.
Similarly, the value of state $s$ under policy $\pi$ is defined as $V^{\pi}(s) = \mathbb{E}\left[R_t|s_t=s\right]$ and is simply the expected return for following policy $\pi$ from state $s$.

In value-based model-free reinforcement learning methods, the action value function is represented using a function approximator, such as a neural network.
Let $Q(s,a;\theta)$ be an approximate action-value function with parameters $\theta$.
The updates to $\theta$ can be derived from a variety of reinforcement learning algorithms.
One example of such an algorithm is Q-learning, which aims to directly approximate the optimal action value function: $Q^*(s,a)\approx Q(s,a;\theta)$.
In one-step Q-learning, the parameters $\theta$ of the action value function $Q(s,a;\theta)$ are learned by iteratively minimizing a sequence of loss functions, where the $i$th loss function defined as
\begin{equation*}
    L_i(\theta_i) = \mathbb{E}\left(r + \gamma \max_{a'}Q(s', a';\theta_{i-1}) - Q(s,a;\theta_i) \right)^2
\end{equation*}
where $s'$ is the state encountered after state $s$.

%
%
%
%
%
%

We refer to the above method as one-step Q-learning because it updates the action value $Q(s,a)$ toward the one-step return $r+\gamma \max_{a'}Q(s',a';\theta)$.
%
One drawback of using one-step methods is that obtaining a reward $r$ only directly affects the value of the state action pair $s,a$ that led to the reward.
The values of other state action pairs are affected only indirectly through the updated value $Q(s,a)$.
This can make the learning process slow since many updates are required the propagate a reward to the relevant preceding states and actions.

%
%

One way of propagating rewards faster is by using $n$-step returns~\citep{watkins1989learning,peng1996msq}. In $n$-step Q-learning, $Q(s,a)$ is updated toward the $n$-step return defined as
%
    $r_t + \gamma r_{t+1} + \cdots + \gamma^{n-1} r_{t+n-1} + \max_a \gamma^n Q(s_{t+n}, a)$.
%
This results in a single reward $r$ directly affecting the values of $n$ preceding state action pairs.
This makes the process of propagating rewards to relevant state-action pairs potentially much more efficient.

%
%

In contrast to value-based methods, policy-based model-free methods directly parameterize the policy $\pi(a|s;\theta)$ and update the parameters $\theta$ by performing, typically approximate, gradient ascent on $\mathbb{E}[R_t]$.
One example of such a method is the REINFORCE family of algorithms due to Williams~\yrcite{Williams1992}.
Standard REINFORCE updates the policy parameters $\theta$ in the direction $\nabla_{\theta}\log\pi(a_t|s_t;\theta) R_t$, which is an unbiased estimate of $\nabla_{\theta} \mathbb{E}[R_t]$.
It is possible to reduce the variance of this estimate while keeping it unbiased by subtracting a learned function of the state $b_t(s_t)$, known as a baseline~\citep{Williams1992}, from the return.
The resulting gradient is
%
$\nabla_{\theta}\log\pi(a_t|s_t;\theta) \left(R_t-b_t(s_t)\right)$.
%

A learned estimate of the value function is commonly used as the baseline $b_t(s_t)\approx V^{\pi}(s_t)$ leading to a much lower variance estimate of the policy gradient.
When an approximate value function is used as the baseline, the quantity $R_t-b_t$ used to scale the policy gradient can be seen as an estimate of the \emph{advantage} of action $a_t$ in state $s_t$, or $A(a_t,s_t)=Q(a_t,s_t)-V(s_t)$, because $R_t$ is an estimate of $Q^{\pi}(a_t,s_t)$ and $b_t$ is an estimate of $V^{\pi}(s_t)$.
This approach can be viewed as an actor-critic architecture where the policy $\pi$ is the actor and the baseline $b_t$ is the critic~\citep{sutton:book,degris2012model}.

%

%
%
%
%
%
%
%
%
%
%
%
%

%
%

